At fixed controller and drive settings, removal of contact and Ib force-feedback pathways reduced stepping in the supported W2L model from 33 left and 28 right cycles to one left and no right cycles. Forward displacement decreased from \qty{1.16}{\meter} to \qty{0.09}{\meter}. These measures cover the \qty{25.499}{\second} scored window after settling within each \qty{30}{\second} simulation. The harness configuration was unchanged, although its measured support contribution differed between conditions. Ia and II pathways remained active. This comparison demonstrates the contribution of the removed pathways at this operating point; it uses the standard muscle-actuated W2L model.
\begin{figure}[t]
\centering
\includegraphics[width=0.85\textwidth]{ablation_pair_metrics_20261026.pdf}
\caption[Contact and force-feedback ablation]{Contact and Ib force-feedback ablation at fixed controller and drive settings in the supported W2L model. Counts and displacement refer to the post-settling scored window. The retained Ia and II pathways distinguish this test from complete deafferentation.}
\label{fig:sns_ablation_pair}
\end{figure}
